\usepackage{parskip}
\usepackage[]{mdframed}
\usepackage{listings}
\definecolor{dkgreen}{rgb}{0,0.6,0}
\definecolor{gray}{rgb}{0.5,0.5,0.5}
\definecolor{mauve}{rgb}{0.58,0,0.82}
\usepackage{tcolorbox}
\usepackage{enumerate}
\usepackage{xparse, pgffor}
\usepackage{pgfplots}
  \pgfplotsset{compat=1.18}
\usepackage{mathrsfs}
% \usepackage[backend=biber, style=numeric]{biblatex}

\lstset{frame=tb,
  language=R,
  aboveskip=3mm,
  belowskip=3mm,
  showstringspaces=false,
  columns=flexible,
  basicstyle={\small\ttfamily},
  numbers=none,
  numberstyle=\tiny\color{gray},
  keywordstyle=\color{blue},
  commentstyle=\color{dkgreen},
  stringstyle=\color{mauve},
  breaklines=true,
  breakatwhitespace=true,
  tabsize=3
}

% Algorithms (algorithmicx)
\usepackage{algorithm}
\usepackage{algpseudocode}
\usepackage{float} % for [H]

% Algorithm (listings-style)
\lstnewenvironment{alg}[1][] {   
  \lstset{ mathescape=true,
      frame=tB,
      numbers=left, 
      numberstyle=\tiny,
      basicstyle=\rmfamily\scriptsize, 
      keywordstyle=\color{black}\bfseries,
      keywords={,procedure, div, for, to, input, output, return, datatype, function, in, if, else, foreach, while, begin, end, }
      numbers=left,
      xleftmargin=.04\textwidth,
      #1
  }
}
{}

% Margins
\setcounter{secnumdepth}{1} 

% Line Spacing in Tables
\newcommand{\tablespace}{\\[1.25mm]}
\newcommand\Tstrut{\rule{0pt}{2.6ex}}         % = `top' strut
\newcommand\tstrut{\rule{0pt}{2.0ex}}         % = `top' strut
\newcommand\Bstrut{\rule[-0.9ex]{0pt}{0pt}}   % = `bottom' strut

% Sets
\renewcommand{\emptyset}{\varnothing}
\newcommand{\R}{\mathbb{R}}
\newcommand{\N}{\mathbb{N}}
\newcommand{\Z}{\mathbb{Z}}
\newcommand{\C}{\mathbb{C}}
\newcommand{\Q}{\mathbb{Q}}
\newcommand{\RQ}{\R\backslash\Q}
\newcommand{\sR}{\mathscr{R}}
\newcommand{\cA}{\mathcal{A}}

% Statistics
\newcommand{\var}{\operatorname{Var}}
\newcommand{\Binom}{\operatorname{Binom}}
\newcommand{\Bias}{\operatorname{Bias}}
\newcommand{\Pois}{\operatorname{Pois}}
\newcommand{\Exp}{\operatorname{Exp}}
\newcommand{\Gammaf}[1]{\Gamma\left(#1\right)}
\newcommand{\E}[1]{\mathbb{E}\left[#1\right]}
\renewcommand{\P}{\mathbb{P}}
\newcommand{\Cov}{\operatorname{Cov}}

% Linear Algebra
\newcommand{\norm}[2]{||#1||_{#2}}
\newcommand{\trace}{\operatorname{trace}}

% Abstract Algebra
\renewcommand{\ker}{\operatorname{Ker}}
\newcommand{\kerphi}{\ker\varphi}
\newcommand{\id}{\operatorname{Id}}
\newcommand{\gen}[1]{\langle #1 \rangle}
\newcommand{\module}[1]{$#1-$\text{module}}
\newcommand{\linear}[1]{$#1-$\text{linear}}
\newcommand{\GL}{\operatorname{GL}}
\newcommand{\SL}{\operatorname{SL}}
\newcommand{\Hom}{\operatorname{Hom}}

% Analysis
\renewcommand{\Re}{\operatorname{Re}}
\renewcommand{\Im}{\operatorname{Im}}
\newcommand{\ovl}[1]{\overline{#1}}
\newcommand{\unl}[1]{\underline{#1}}
\newcommand{\vepsi}{\varepsilon}
\newcommand{\diam}{\operatorname{diam}}
\newcommand{\vvf}{\textbf{f}}

% Topology
\newcommand{\Cl}{\operatorname{Cl}}

% Calculus
\newcommand{\del}{\partial}

% Computer Science
\newcommand{\arr}{\texttt{[]}}
\newcommand{\true}{\texttt{True}}
\newcommand{\false}{\texttt{False}}
\newcommand{\dpt}{\texttt{dp}} % do not make this \dp, apparently that is used as an internal command in LaTeX and it causes errors when we try to redefine it, so we will just use \dpt instead for our dynamic programming table! insane!

% Misc
\newcommand{\ceil}[1]{\left\lceil #1 \right\rceil}
\newcommand{\floor}[1]{\left\lfloor #1 \right\rfloor}
\newcommand{\notexists}{\nexists}
\newcommand{\im}{\operatorname{im}}
\renewcommand{\frak}[1]{\mathfrak{#1}}
\NewDocumentCommand{\quadn}{O{1}}{\foreach \n in {1,...,#1}{\quad} }
\newcommand{\Proof}{\textbf{Proof: }}
\newcommand{\hint}{$\mathbf{Hint\!\!:}\;\;$}
\newcommand{\TODO}[1]{%
  \begin{tcolorbox}[colback=yellow!20, colframe=red, title=TODO]
  #1
  \end{tcolorbox}
}
\newcommand{\bmid}{\;\bigg|\;}
\DeclareMathOperator*{\argmax}{arg\,max}
\DeclareMathOperator*{\argmin}{arg\,min}

% Syntactic Proofs
\newenvironment{synproof}
    {\begin{enumerate}[$t_1: \quad$]} % start enumerate with custom labels
    {\end{enumerate}}            % end enumerate
\newcommand{\step}[2]{\item {#1} \hfill \text{#2}}
